\documentclass{situationist-preprint}

\title{Indexicality and Presentness}
\msid{TSM-15-b}
\msfamily{TSM-15}
\msdiscipline{Time, Space and Modality}
\mstype{Companion manuscript}
\mskeywords{presentness; indexicals; self-locating belief; A-theory; essential indexical}

\begin{document}
\maketitle

\begin{abstract}
TSM-15-a argues that the present has no metaphysical privilege: it is the time at which a thought or utterance occurs. The A-theorist replies that learning that it is now noon is learning something that no tenseless description conveys, and that what is learned is a fact about which time is present. We argue that the information is self-locating. A subject who knows the entire tenseless history of the world can still be ignorant of when she is, as she can be ignorant of where she is; the information she lacks is a centring of what she knows, not a further fact. The essential indexical shows the irreducibility of the first-person perspective, not the privilege of a time.
\end{abstract}

\section{The argument from tensed information}

\citet{Prior1959} and \citet{Perry1979} showed that tensed beliefs cannot be replaced by tenseless ones without loss. Someone who believes that the meeting starts at noon on 3 May, and does not believe that it is now noon on 3 May, does not get up and go; once she believes that it is now noon, she does. The A-theorist concludes that \enquote{it is now noon} reports a fact, the presentness of noon, that the tenseless belief omits. If so, Assumption~2.3 of TSM-15-a is false: there is a monadic presentness, and one time has it.

\section{Self-locating information}

\citet{Lewis1979} proposed that the objects of belief are properties, or sets of centred worlds: pairs of a world and a subject-at-a-time within it. To believe that it is now noon is to self-ascribe the property of being located at noon. EPI-18-c applies this to Sleeping Beauty: she knows the history of her world, apart from the coin, and is uncertain only where in it she is; her learning is conditionalization on self-locating information.

\begin{proposition}\label{p:centre}
A subject who knows every tenseless truth may still lack the belief that it is now $t$; the belief she lacks is self-locating, and its truth requires no feature of $t$ other than her thought's occurring at $t$.
\end{proposition}

\begin{proof}
Tenseless knowledge fixes a set of worlds; it does not fix a centre within the actual world, since the subject's thoughts occur at many times and the tenseless truths are the same at all of them. The belief that it is now $t$ narrows the centres to those at $t$. It is true, as a centred belief, exactly when the thought occurs at $t$, which by Definition~2.1 of TSM-15-a is for $t$ to be present relative to it.
\end{proof}

The proposition explains the data of the argument without the A-theorist's fact. The irreducibility is real: no uncentred belief does the work of the centred one. But it is irreducibility of the perspective of the believer, and it holds of \enquote{here} and \enquote{I} as much as of \enquote{now} \citep{Perry1979}. The amnesiac in the library who knows every fact in the books may not know that he is Rudolf Lingens; nobody infers a monadic property of being me that one person has and the others lack. The parallel conclusion for times is Theorem~3.1 of TSM-15-a: no time is privileged by presentness.

\begin{objection}
Spatial self-location does not change with no motion of mine, but temporal self-location always does: what is now changes whether I act or not. Presentness moves, and hereness does not.
\end{objection}

\begin{reply}
A subject who persists through time has thoughts at many times, and each is centred on its own time; there is no single centre that moves. A subject who moves through space likewise has thoughts at many places. The difference is that persisting subjects cannot avoid occupying successive times, which is a fact about the ordering of their thoughts in the B-series, not about a property that passes from time to time.
\end{reply}

\begin{objection}
Centred worlds presuppose times; the A-theorist may say that the centres are A-determinations.
\end{objection}

\begin{reply}
Centres are times, places and subjects, all of which the tenseless view has. The A-theorist must add that one centre is privileged absolutely, and nothing in the explanation of action or learning uses that addition.
\end{reply}

\section{Conclusion}

The information conveyed by \enquote{it is now noon} is self-locating: it locates the thinker in a history she may already know. Indexicality explains why it cannot be stated tenselessly, and explains it in the same way for times, places and persons. Presentness is the relation of a time to the thoughts occurring at it, and the present is not special.

\begin{thebibliography}{99}\small
\bibitem[Lewis(1979)]{Lewis1979} Lewis, D. (1979). Attitudes de dicto and de se. \emph{The Philosophical Review}, 88(4), 513--543.
\bibitem[Perry(1979)]{Perry1979} Perry, J. (1979). The problem of the essential indexical. \emph{No\^us}, 13(1), 3--21.
\bibitem[Prior(1959)]{Prior1959} Prior, A.~N. (1959). Thank goodness that's over. \emph{Philosophy}, 34(128), 12--17.
\end{thebibliography}

\end{document}
